
\sourcenote{Genre: core rule (tier 1) --- an ordinary mathematical proof at
Putnam-exam-solution detail. Every hypothesis of the rule is used below and
remarked on, with the counterexample its removal admits where one exists;
the formal-anchor table maps each step to the Lean declaration that
carries it.}

\textbf{Setting and claim.}\enspace Fix a hypernet $H$. A valuation
$\sigma$ assigns each wire a rational tensor of its declared shape; a
sequent $\Phi \vdash \Psi$ is semantically valid when every admissible
$\sigma$ satisfying $\Phi$ satisfies $\Psi$, with comparison atoms read
per-entry. A region atom $x \in \mathcal{R}(\Delta, S)$, for
$\Delta = \mathrm{conv}(V_0, \dots, V_n)$ an $n$-simplex and
$S \subsetneq \{0, \dots, n\}$ a proper set of \emph{extended} vertices,
holds at $\sigma$ when the evaluated point admits barycentric weights
$\lambda_v$ with $\sum_v \lambda_v = 1$, $\lambda_v \geq 0$ at core
vertices $v \notin S$, and $\lambda_v \leq 0$ at $v \in S$ --- the closed
simplex when $S = \emptyset$, and the unbounded set extending beyond the
faces opposite $S$ otherwise. The rule's payload supplies the region
premise $\Phi \vdash x \in \mathcal{R}(\Delta, S)$, a polynomial term $p$
of degree $\leq d$, a predicate $P$ from the accepted grammar (a finite
conjunction of scalar hole-atoms $s \cdot \square \bowtie c$, $s > 0$,
$\bowtie \in \{\leq, \geq, =\}$), the checked wire tie \texttt{polyTie},
the predicate-shape clause ($P$ conic when $S \neq \emptyset$; convex
suffices when $S = \emptyset$), and the checked coefficient condition:
every forward-difference Bernstein coefficient of $p$ over
$\mathcal{R}(\Delta, S)$, recomputed exactly, satisfies $P$. The claim
(\texttt{sound\_bernCone}): the composite sequent
$\Phi \vdash P(p(x))$ is then semantically valid.

\begin{blemma}[Bernstein form on the closed simplex]\label{lem:bc-bern}
Let $p$ be a polynomial of degree $\leq d$ on $\R^n$ with degree-$d$
Bernstein coefficients $(b_\alpha)_{|\alpha| = d}$ on $\Delta$, so that
$p(x) = \sum_{|\alpha| = d} b_\alpha\, B^d_\alpha(\lambda)$ with
$B^d_\alpha(\lambda) = \binom{d}{\alpha} \lambda^\alpha$ in the
barycentric coordinates $\lambda$ of $x$. For $x \in \Delta$ (all
$\lambda_v \geq 0$),
\[
B^d_\alpha(\lambda) \geq 0
\qquad\text{and}\qquad
\sum_{|\alpha| = d} B^d_\alpha(\lambda)
  = \Bigl(\sum_v \lambda_v\Bigr)^{\!d} = 1,
\]
so $p(x)$ is a \emph{convex combination} of the coefficients $b_\alpha$.
\end{blemma}
\begin{proof}
Nonnegativity is termwise ($\lambda^\alpha \geq 0$); the partition of
unity is the multinomial theorem applied to
$(\lambda_0 + \dots + \lambda_n)^d = 1^d$. The representation itself is
the classical change to the Bernstein basis
\cite{bernstein1912,farouki2012}; the checker \emph{computes} the
$b_\alpha$ from $p$'s canonical monomial form by exact forward
differencing over $\Q$, so the representation is an identity of exact
rational data, not an approximation. That Bernstein coefficients
enclose the range of $p$ on $\Delta$ --- the reading of this lemma used
below --- is the sound (easy) direction of \cite{cargoshisha1966}.
\end{proof}

\begin{blemma}[accepted predicates are convex]\label{lem:bc-convex}
Let $P$ be accepted by the predicate classifier. Then its satisfaction
set $\{ v \in \Q : \mathrm{satisfiesPredQ}(P, v) \}$ is convex.
\end{blemma}
\begin{proof}
Each accepted atom orients (possibly by mirroring) to
$s \cdot v \bowtie c$ with $s > 0$ and
$\bowtie \in \{\leq, \geq, =\}$; its satisfaction set is a closed ray or
a point --- an interval either way --- and division by the positive $s$
does not disturb this. A conjunction intersects finitely many intervals.
Strict comparators and $\neq$ are \emph{refused} by the classifier ---
closed-world conservatism that can lose completeness but never
soundness; the refusal is a checker verdict, not a soundness case.
\end{proof}

\begin{blemma}[the univariate ray identity]\label{lem:bc-star}
Let $d \geq 0$ and let $b_0, \dots, b_d$ be the degree-$d$ Bernstein
coefficients of $p$ on the segment $[V_a, V_b]$, indexed by the
multiplicity of $V_b$. Let $x$ lie on the extension ray beyond $V_b$:
in barycentric coordinates $x = \lambda_a V_a + \lambda_b V_b$ with
$\lambda_a = -u \leq 0$, $\lambda_b = 1 + u$, $u \geq 0$. Then, with
$(\Delta b)_i := b_{i+1} - b_i$ the forward difference,
\[
p(x) \;=\; \sum_{j=0}^{d} \binom{d}{j}\, u^{j}\, (\Delta^{j} b)_{d-j},
\tag{$\star$}
\]
a combination of the certified table values $(\Delta^{j} b)_{d-j}$ with
\emph{nonnegative} weights $\binom{d}{j} u^{j}$, the $j = 0$ weight
being $1$.
\end{blemma}
\begin{proof}
Substitute $\lambda_a = -u$, $\lambda_b = 1 + u$ into the Bernstein
form and expand $(1+u)^i$ binomially:
\[
p(x) = \sum_{i=0}^{d} \binom{d}{i} b_i\, (-u)^{d-i} (1+u)^{i}
     = \sum_{i=0}^{d} \sum_{k=0}^{i} \binom{d}{i} \binom{i}{k}
       (-1)^{d-i}\, b_i\, u^{d-i+k}.
\]
Collect the coefficient of $u^{j}$: setting $i = d - j + m$ with
$m = 0, \dots, j$, it is
$\sum_{m=0}^{j} \binom{d}{d-j+m} \binom{d-j+m}{m} (-1)^{j-m}
b_{d-j+m}$. The product of binomials telescopes,
\[
\binom{d}{d-j+m} \binom{d-j+m}{m}
  = \frac{d!}{(j-m)!\, m!\, (d-j)!}
  = \binom{d}{j} \binom{j}{m},
\]
so the coefficient of $u^{j}$ is
$\binom{d}{j} \sum_{m=0}^{j} \binom{j}{m} (-1)^{j-m} b_{d-j+m}
 = \binom{d}{j} (\Delta^{j} b)_{d-j}$, which is $(\star)$. (Degree-$1$
sanity check: $p(x) = b_1 + u\,(b_1 - b_0)$, matching
$(1+u)\,b_1 - u\,b_0$ directly.) The values $(\Delta^{j} b)_{d-j}$ are
exactly the rows the checker's forward-difference table certifies on
this fragment: the $j$-th difference of the coefficient sequence,
anchored at the core vertex.
\end{proof}

\begin{blemma}[conic transfer]\label{lem:bc-conic}
Let $P$ be accepted with the conic classification (every atom's
constant is $0$). Then its satisfaction set contains $0$, and is closed
under nonnegative scaling and under finite nonnegative combinations.
\end{blemma}
\begin{proof}
Each conic atom reads $s \cdot v \bowtie 0$, satisfied at $v = 0$ and
preserved under $v \mapsto t v$ for $t \geq 0$; conjunction preserves
both facts. Closure under addition follows from convexity
(Lemma~\ref{lem:bc-convex}) plus scaling:
$v + w = 2 \cdot \tfrac{1}{2}(v + w)$; induction extends to finite
nonnegative combinations.
\end{proof}

\begin{bthm}[soundness of \rulename{BernCone}]\label{thm:bc-sound}
Under the payload's checked hypotheses, if the region premise is
semantically valid, then so is $\Phi \vdash P(p(x))$.
\end{bthm}
\begin{proof}
Fix an admissible $\sigma \models \Phi$; the region premise yields
barycentric weights for the evaluated point of $x$ as in the Setting.
The conclusion formula is, by the checked \texttt{postApply}, the
predicate $P$ with the polynomial term substituted for its hole, so it
suffices to establish $\mathrm{satisfiesPredQ}(P, \cdot)$ at the
evaluated value(s) of $p$ --- per-entry, when $p$ is tensor-shaped.

\emph{Step 1 (the polynomial reads one point).} The wire tie
\texttt{polyTie} says $p$ is closed, or mentions exactly the region
wire $w$ with $x = w$. In the second case the region premise pins
$\sigma$'s read of $w$ to a point of $\mathcal{R}(\Delta, S)$; in the
first, no wire is read at all and $p$'s value is its closed constant.

\emph{Step 2 (rational weights).} The model's barycentric weights are
\emph{real}; the kernel computes over $\Q$. The weights solve a linear
system with rational data (the vertex coordinates and the evaluated
point), so a nonnegative real solution can be traded for a nonnegative
rational one, and the exact kernel consumes the rational weights. This
step is invisible in informal treatments and essential in the
mechanization.

\emph{Step 3 (case $S = \emptyset$: the convex kernel).} By
Lemma~\ref{lem:bc-bern}, $p$'s value is a convex combination of the
certified coefficients; every coefficient satisfies $P$ (the checked
table sweep); by Lemma~\ref{lem:bc-convex} satisfaction transfers to
the combination. On this case the checker admits both translation
routes --- the idx-free single-wire form and the multivariate
projection-indexed form with all indices in range --- and the per-entry
evaluation lemma carries the scalar argument to every entry of a
tensor-shaped value.

\emph{Step 4 (case $S \neq \emptyset$).} The checker's scope guard
confines accepted certificates to two shapes. \emph{(4a: constant
canonical form, any dimension.)} If $p$'s canonical polynomial has no
variables, its value everywhere is the constant --- the zero-order
forward-difference row is that constant, certified directly.
\emph{(4b: single-wire polynomial over a 1-dimensional region.)} A
proper nonempty extension set on a 1-simplex extends exactly one of the
two vertices, and $\mathcal{R}(\Delta, S)$ is the closed ray from the
core vertex; Lemma~\ref{lem:bc-star} writes $p$'s value as a
nonnegative combination of the certified table values with the $j = 0$
weight equal to $1$, and Lemma~\ref{lem:bc-conic} transfers
satisfaction through the combination. Everything else is
\emph{refused} by the guard, so no further case exists for soundness
to cover.

\emph{Step 5 (read off).} In each case
$\mathrm{satisfiesPredQ}(P, \cdot)$ holds at the evaluated value; the
substitution reading turns this into
$\sigma \models P(p(x))$.
\end{proof}

\begin{bremark}[why \texttt{polyTie} is not decoration]
\label{rem:bc-polytie}
Nothing else ties the checked polynomial to the wire the premise
constrains. Without the field, take the satisfiable premise
$w_0 \in \mathcal{R}(\mathrm{conv}(1,2), \emptyset)$ and the polynomial
$p = w_1$ over a \emph{different} wire: the translation remaps the
single-variable canonical form onto the region coordinate, the
certified coefficients are $\{1, 2\}$ --- both satisfying
$\square \geq 1$ --- and the rule would conclude
$\Phi \vdash w_1 \geq 1$, false at any valuation reading
$\sigma(w_1) = 0$. Closed polynomials need no tie: their value reads no
wire.
\end{bremark}

\begin{bremark}[why conic, not merely convex, on rays]
\label{rem:bc-conic-needed}
The predicate-shape clause is two-sided because the two region shapes
transfer differently: convex combinations
(Lemma~\ref{lem:bc-bern}) on the closed simplex, conic combinations
with unbounded weights ($\star$) on a ray. A convex-but-not-conic
predicate fails on the ray, by hand: take $d = 1$,
$\Delta = \mathrm{conv}(0, 1)$ with the vertex at $0$ extended (the
region is the ray $[1, \infty)$), $p(t) = 1 - t$, and
$P = (\square \geq -1)$. The certified table is
$\{\,b_1,\ (\Delta b)_0\,\} = \{\,0,\ -1\,\}$, every value satisfying
$P$; yet at $u = 2$ (the point $t = 3$), $(\star)$ gives
$p = b_1 + 2(b_1 - b_0) = -2 < -1$. The classifier refuses exactly
this: an atom with nonzero constant classifies non-conic, and the
$S \neq \emptyset$ clause then rejects the certificate. Conversely, at
$S = \emptyset$ demanding conic would be needless strength --- the rule
carries the license (any accepted convex predicate) in the premise
itself rather than in surrounding prose.
\end{bremark}

\begin{bremark}[the scope guard is part of the theorem]
\label{rem:bc-guard}
The checker refuses non-constant idx-free polynomials over regions of
dimension $\neq 1$, and multivariate projection-indexed polynomials
whenever $S \neq \emptyset$. Soundness is stated and proven for every
certificate the checker \emph{accepts}; the guard makes the unproven
territory unrepresentable rather than discouraged. The fully general
multivariate ray identity (the arbitrary-shape analogue of
Lemma~\ref{lem:bc-star}) remains open in the development and becomes
essential only if the guard lifts --- the two must move together.
\end{bremark}

\begin{bremark}[relation to the literature]
The rule is the certificate form of Bernstein-coefficient range
enclosure: coefficients bound the range on the simplex
\cite{cargoshisha1966}, enclosures converge under degree elevation and
subdivision \cite{garloff1986} --- which is why \rulename{Subdivide} is
this rule's completeness partner --- and the basis itself goes back to
\cite{bernstein1912}, surveyed in \cite{farouki2012}. With polynomial
atoms and positivity certificates the rule sits beside the
Krivine--Handelman family \cite{krivine1964,handelman1988}: as
everywhere in this calculus, only the trivial (sound) direction is
claimed or needed --- no Positivstellensatz completeness enters the
kernel. The forward-difference computation of the coefficient table is
the classical finite-difference reading of the Bernstein form; the
subdivision algorithms of \cite{decasteljau1959,laneriesenfeld1980}
operate on the same data one rule over.
\end{bremark}

\begin{anchors}
Setting: region atoms as masked barycentric membership & \texttt{MaskRegion}; \texttt{Formula.holds} (\texttt{inRegion} case)\\
Lemma \ref{lem:bc-bern} + Step 3 kernel & the $S = \emptyset$ convex kernel consumed by \texttt{sound\_bernCone\_empty}; per-entry carry \texttt{Term.evalQ\_brRead}\\
Step 2 rational weights & \texttt{maskRegion\_rat\_weights}; \texttt{ConvexRat.exists\_nonneg\_rat\_solution}\\
Lemma \ref{lem:bc-convex} & \texttt{classifyPred}, \texttt{satisfiesPredQ\_convex}\\
Lemma \ref{lem:bc-star} (univariate $(\star)$) & \texttt{QArith.fdBernstein\_sound\_dim1}\\
Step 4a constant corner & \texttt{QArith.fdBernstein\_sound\_const}\\
Lemma \ref{lem:bc-conic} & \texttt{satisfiesPredQ\_conic\_zero}, \texttt{satisfiesPredQ\_conic\_scale}\\
coefficient recomputation (the checked table) & \texttt{bernConeCoeffCheck}, \texttt{allForwardDiffBernCoeffsSatisfyP}, \texttt{QArith.fdBernstein}, \texttt{QArith.gridAll}\\
coordinate translation + index gate & \texttt{Term.asCoordPoly?}, \texttt{Term.asCoordPolyIdx?}, \texttt{Term.idxBounded}\\
Step 1 wire tie & payload field \texttt{polyTie}\\
Step 5 conclusion reading & \texttt{applyPredicateToTerm}, \texttt{holds\_substHole\_of\_satisfies}\\
Theorem \ref{thm:bc-sound} & \texttt{sound\_bernCone} (case split onto \texttt{sound\_bernCone\_empty}); wired as core \texttt{RuleData.bernCone} into \texttt{ProofObject.sound}\\
open residue (Remark \ref{rem:bc-guard}) & \texttt{ConeLemma.StarIdentity} (general shapes; not consulted by the current guard)\\
\end{anchors}
